\documentclass[12pt,reqno]{amsart}

\usepackage{amsmath,amssymb,amsthm,mathtools}
\usepackage[utf8]{inputenc}
\usepackage[T1]{fontenc}
\usepackage[hidelinks]{hyperref}
\usepackage{booktabs}
\usepackage[nameinlink,noabbrev]{cleveref}
\setlength{\emergencystretch}{2em}

\theoremstyle{plain}
\newtheorem{theorem}{Theorem}[section]
\newtheorem{proposition}[theorem]{Proposition}
\newtheorem{lemma}[theorem]{Lemma}
\newtheorem{corollary}[theorem]{Corollary}
\theoremstyle{remark}
\newtheorem{remark}[theorem]{Remark}
\theoremstyle{definition}
\newtheorem{definition}[theorem]{Definition}
\newtheorem{example}[theorem]{Example}
\newtheorem*{theoremA}{Theorem A}
\newtheorem*{theoremB}{Theorem B}
\newtheorem*{theoremC}{Theorem C}

\numberwithin{equation}{section}
\newcommand{\ZZ}{\mathbb{Z}}
\newcommand{\RR}{\mathbb{R}}
\newcommand{\EE}{\mathbb{E}}
\newcommand{\card}[1]{\lvert #1 \rvert}

\title[Brocot condensation and critical renewal]{Brocot Condensation and\
Critical Fibonacci Renewal}

\author{Haobo Ma}
\address{AELF PTE LTD., \#14-02, Marina Bay Financial Centre Tower 1,
8 Marina Blvd, Singapore 018981, Singapore}
\email{auric@aelf.io}

\author{Wenlin Zhang}
\address{National University of Singapore, Singapore}
\email{e1327962@u.nus.edu}

\hypersetup{
  pdftitle={Brocot Condensation and Critical Fibonacci Renewal},
  pdfauthor={Haobo Ma and Wenlin Zhang},
  pdfsubject={Denominator-weighted Brocot fractions and critical Fibonacci renewal},
  pdfkeywords={Brocot fractions, continued fractions, continuants, condensation, Fibonacci partition function, renewal theory}
}
\date{}
\subjclass[2020]{11A55, 11B39, 60F05, 60K05}
\keywords{Brocot fractions, continued fractions, continuants, condensation,
Fibonacci partition function, renewal theory, stable laws}

\begin{document}

\begin{abstract}
Fix \(s>2\), and weight the canonical regular continued fractions
\([0;a_1,\ldots,a_r]\) with \(a_r\ge2\) and
\(a_1+\cdots+a_r=n\) proportionally to the inverse \(s\)-th power of their
denominators.  We prove a total-variation condensation theorem.  With
probability tending to one there is a unique partial quotient larger than
\(n/2\), and the words on its two sides converge jointly to independent
finite words with inverse-continuant laws.  This gives the exact limiting
defect and location laws and an asymptotic factorization of the denominator.

The normalization is the fixed-digit-sum continuant sum studied by
Dushistova after the Farey-tree work of Moshchevitin and Zhigljavsky.  We
prove
\[
 Z_n(s)\sim
 2\left\{\frac{\zeta(s-1)}{\zeta(s)}\right\}^{\!2}n^{-s}
\]
without a tail hypothesis.  An exact context calculation identifies an
endpoint overcount in Dushistova's argument and changes its printed leading
coefficient from \(R_s+2R_s^2\) to \(2R_s^2\); no assertion is made about the
remaining coefficients of that expansion.

At the arithmetic critical point
\(\zeta(\sigma_0-1)/\zeta(\sigma_0)=2\), the scalar estimate gives an exact
regularly varying cost tail with context constant \(8\).  Combining it with
the free-generator coordinates of Weinstein yields a weighted renewal
identity for the classical Fibonacci partition function, a spectrally
positive stable domain of attraction, and the noninteger
\(m^{3-\sigma_0}\) correction to its critical one-layer partition sum.
\end{abstract}

\maketitle

\section{Introduction}\label{sec:introduction}

The Stern--Brocot and Farey constructions organize rational numbers by
continued-fraction complexity.  A natural statistical question is what a
typical rational looks like when a fixed complexity layer is biased toward
small denominators.  In the model considered here, the complexity of the
canonical regular continued fraction
\[
 x=[0;a_1,\ldots,a_r],\qquad a_i\ge1,\quad a_r\ge2,
\]
is the digit sum \(a_1+\cdots+a_r=n\), and the weight is
\(\operatorname{den}(x)^{-s}\) for a fixed \(s>2\).  The continuant in the
denominator is nonmultiplicative across the digits.  Nevertheless, a single
partial quotient absorbs all but a tight amount of the digit sum, while the
two finite contexts on its sides converge to independent arithmetic laws.

This condensation statement is stronger than the scalar asymptotic for the
normalizing sum.  It identifies the limiting context words in total
variation, the defect of the largest partial quotient, its position in the
continued fraction, the numbers of digits on either side, and the leading
factorization of the sampled denominator.  The proof also gives the scalar
asymptotic with its exact constant.  That scalar consequence becomes the
arithmetic input for a different but closely linked question: the critical
renewal structure of the classical Fibonacci partition function.

The link is supplied by Weinstein's free generators for Fibonacci
partitions \cite{Weinstein2016}.  Their letters are reduced fractions in
\((0,1)\); a letter \(p/q\) carries multiplicative weight \(q\) and an odd
cost read from its negative continued fraction.  The cost-\(2d+1\) letter
sum is termwise identical to the regular-continuant sum at digit sum
\(d+1\).  Thus the same scalar estimate that normalizes the condensation
law determines the critical generator-cost tail.  At the unique
\(\sigma_0>2\) satisfying
\[
 \frac{\zeta(\sigma_0-1)}{\zeta(\sigma_0)}=2,
\]
the generator letters form a probability distribution with a regularly
varying tail of index \(1-\sigma_0\).  Its mean is finite and its variance is
infinite.  A second-order lattice renewal theorem then produces a fractional
power correction to the critical Fibonacci partition sum, and the same tail
places the cost in the domain of attraction of a spectrally positive stable
law.

\subsection{Main results}

For a finite word \(\mathbf w\), let \(\mathcal K(\mathbf w)\) be its
regular continuant and let \(|\mathbf w|_1\) be its digit sum.  Write
\(\mathcal W_{\rm L}\) for all finite positive words, including the empty
word, and \(\mathcal W_{\rm R}\) for the canonical regular words, again
including the empty word.  Put
\[
 \rho_s=\frac{\zeta(s-1)}{\zeta(s)}.
\]
The precise definitions occur in \Cref{sec:brocot-condensation}.  The first
result is summarized as follows.

\begin{theoremA}[Total-variation condensation]
For fixed \(s>2\), sample a canonical continued fraction of digit sum \(n\)
with probability proportional to the inverse \(s\)-th power of its
denominator.  With probability tending to one it has a unique digit larger
than \(n/2\).  If \(U_n,V_n\) are the words to the left and right of that
digit, then
\[
 (U_n,V_n)\xrightarrow{\rm TV}(U,V),
\]
where \(U,V\) are independent and
\[
 \Pr\{U=\mathbf u\}=\frac{\mathcal K(\mathbf u)^{-s}}{2\rho_s},
 \qquad
 \Pr\{V=\mathbf v\}=\frac{\mathcal K(\mathbf v)^{-s}}{\rho_s}.
\]
The defect, the location data, and the normalized denominator converge as
stated in \Cref{thm:brocot-condensation}.
\end{theoremA}

The asymmetry between the two context laws is caused only by the canonical
terminal-digit convention.  A regular rational has two positive
continued-fraction expansions, whereas the canonical form retains the one
whose final digit is at least two.  The empty word and the one-letter word
\((1)\) are endpoint exceptions.  Accounting for those exceptions is also
what resolves the constant in the earlier scalar calculation.

Let
\[
 Z_n(s)=\sum_{\substack{a_1+\cdots+a_r=n\\a_i\ge1,\ a_r\ge2}}
       \mathcal K(a_1,\ldots,a_r)^{-s}.
\]
Dushistova's notation is
\(\sigma_\beta(n)=Z_n(2\beta)\).  The earlier paper obtains a fuller
expansion whose displayed leading coefficient is
\(R_s+2R_s^2\), where \(R_s=\rho_s\)
\cite{Dushistova2007}.  The exact endpoint decomposition gives the following
result.

\begin{theoremB}[Context constant and denominator layer]
For every fixed \(s>2\),
\[
 Z_n(s)\sim2\rho_s^2n^{-s}.
\]
Under the termwise change from regular to negative continued fractions, the
cost-\(2d+1\) denominator layer satisfies
\[
 b_{2d+1}(s)=Z_{d+1}(s)\sim2\rho_s^2d^{-s}.
\]
The constant \(2R_s^2\) corrects the leading coefficient in Dushistova's
expansion.  No conclusion is drawn here about its other coefficients.
\end{theoremB}

The proof is deliberately local.  The left context of digit sum zero
contributes \(R_s\), the context of digit sum one contributes another
\(R_s\), and the contexts of larger digit sum contribute
\(2(R_s-1)R_s\).  Their sum is \(2R_s^2\).  In the printed convolution the
condition excluding the empty left context is lost after the \(u=1\) term is
separated; because the remaining convolution is doubled, that empty context
is counted twice.  Restoring one subtraction removes precisely the excess
\(R_s\).  This is not a change of continued-fraction convention or exponent.

For the probabilistic theorem, the central estimate is a summable
two-context domination.  Once a digit \(a>n/2\) is singled out,
multilinearity of the continuant gives
\[
 \mathcal K(\mathbf u,a,\mathbf v)
 =a\mathcal K(\mathbf u)\mathcal K(\mathbf v)+O_{\mathbf u,\mathbf v}(1).
\]
The inverse-continuant mass of the two contexts is summable for \(s>2\).
Words with no digit above \(n/2\) are negligible: a greedy block
decomposition handles words made only of moderate digits, and a direct cut
handles the cases with one or at least two larger digits.  Dominated
convergence gives the product context law, and the discrete Scheff\'e
argument upgrades pointwise convergence to total variation.

The second part of the paper transfers only the scalar layer estimate, not
the total-variation theorem.  Let \(F_0=0,F_1=1\), and let \(R(N)\) count
representations of \(N\) as a sum of distinct members of
\(F_2,F_3,\ldots\).  On the standard Fibonacci layer
\[
 \mathcal I_m=[F_{m+1}-1,F_{m+2}-1),
\]
put
\[
 Z_m^R(t)=\sum_{N\in\mathcal I_m}R(N)^t.
\]
The free-generator parameterization yields a renewal sequence \(u_j(s)\)
whose exact layer identity is
\[
 Z_m^R(-s)=2\sum_{j=0}^{m-1}u_j(s)+u_m(s)-1.
\]
This identity and its proof are given in
\Cref{prop:weighted-generator-renewal}.

At \(s=\sigma_0\), let \(C\) be the generator-letter cost.  Set
\(\alpha=\sigma_0-1\in(1,2)\), and let \(\mu_C=\EE C\).  The arithmetic
constant from the two context masses is \(b_C=8\), while
\[
 K_C=\frac{2^{\sigma_0-1}}{\sigma_0-1}b_C.
\]

\begin{theoremC}[Critical renewal and stable law]
The critical letter law satisfies
\[
 \Pr\{C>x\}\sim K_Cx^{1-\sigma_0}.
\]
The associated renewal sequence has the second-order expansion
\[
 u_j(\sigma_0)=\frac1{\mu_C}
 +\frac{K_C}{\mu_C^2(\sigma_0-2)}j^{2-\sigma_0}
 +o(j^{2-\sigma_0}),
\]
and hence
\[
 Z_m^R(-\sigma_0)=\frac{2m}{\mu_C}
 +\frac{2K_C}{\mu_C^2(\sigma_0-2)(3-\sigma_0)}
  m^{3-\sigma_0}+o(m^{3-\sigma_0}).
\]
Moreover, centered sums of independent copies of \(C\), normalized by a
sequence \(a_n\) with \(n\Pr\{C>a_n\}\to1\), converge to the centered
spectrally positive \(\alpha\)-stable law specified in
\Cref{thm:arithmetic-critical-stable-renewal}.  One may take explicitly
\[
 a_n=\left(\frac{2^{\sigma_0+2}}{\sigma_0-1}n\right)^{1/(\sigma_0-1)}.
\]
\end{theoremC}

The noninteger power is a finite-size signature of the arithmetic critical
tail.  Regular variation is not assumed in the renewal step: it is proved
first from the denominator-layer asymptotic.  The second-order renewal
theorem of Omey and Van Gulck \cite{OmeyVanGulck2015} then applies because
the law is aperiodic, supported on the positive integers, and has finite
mean.  The stable conclusion follows from the standard one-sided
domain-of-attraction criterion \cite{Feller1971}.

\subsection{History and contribution boundary}

Moshchevitin and Zhigljavsky studied entropy sums associated with partitions
of the unit interval generated by the Farey tree
\cite{MoshchevitinZhigljavsky2004}.  Dushistova developed precise
fixed-digit-sum denominator estimates for Brocot sequences
\cite{Dushistova2007}.  The present scalar sum is exactly her local sum after
the exponent change \(s=2\beta\), and the relation is termwise rather than
asymptotic.  The contribution here is the joint total-variation structure,
the endpoint calculation of the leading constant, and the shorter scalar
proof that does not assume an auxiliary tail estimate.  The correction is
confined to the leading coefficient: the arguments here neither verify nor
reject the later coefficients in the published expansion.

Condensation under a large-sum conditioning is familiar for independent
heavy-tailed variables and for probabilistic combinatorial structures.  In
particular, total-variation one-big-jump results and Gibbs-partition
condensation provide useful comparisons
\cite{ArmendarizLoulakis2011,Stufler2020}.  They do not imply the theorem
here.  A continuant weight does not factor into weights of its digits, so the
model is not an independent allocation or a multiplicative Gibbs partition.
The proof instead exploits the exact two-sided continuant inequality and the
summability of all finite contexts.  Accordingly, the result is presented as
a structural theorem for this nonmultiplicative arithmetic family, not as a
general condensation principle.

Fibonacci representations have a long history, beginning with the canonical
representation theorem and continuing through exact formulas and extremal
questions \cite{Zeckendorf1972,Carlitz1968,Berstel2001}.  Weinstein's orbit
and generator theory gives the coordinate system used here
\cite{Weinstein2016}.  We do not reproduce that classification.  We state
the precise pieces of it needed for the weighted layer count, check the
Fibonacci indexing and endpoint convention, and derive the renewal identity.
The critical tail constant, the second-order renewal term, and the stable
limit are then consequences of the continuant calculation in this paper.

A companion manuscript by the same authors, \emph{Finite-window Zeckendorf
thermodynamics}, studies an exact two-layer finite-window correspondence,
critical coexistence, and large deviations across the pressure corner.  No
result from that manuscript is used here.  Conversely, its leading critical
renewal law requires only the finite-mean renewal theorem; the exact tail and
fractional correction proved here are not prerequisites for its principal
finite-window results.

\subsection{Organization}

\Cref{sec:brocot-condensation} defines the denominator-weighted Brocot
model, proves the context estimates and the total-variation theorem, and
derives the corrected denominator-layer asymptotic.
\Cref{sec:fibonacci-renewal} introduces the classical Fibonacci partition
function, the minimal Weinstein generator coordinates, and the exact
one-layer weighted renewal identity.  The arithmetic critical tail,
second-order renewal expansion, stable law, and Fibonacci finite-size
consequence are proved in \Cref{sec:arithmetic-critical}.  The final section
records the scope of the two mechanisms and the open directions left by the
argument.

\section{Condensation for denominator-weighted Brocot fractions}
\label{sec:brocot-condensation}

For a finite word \(\mathbf w=(w_1,\ldots,w_k)\) of positive integers,
let \(\mathcal K(\mathbf w)\) denote its regular continuant, with
\[
 \mathcal K(\varnothing)=1,\qquad
 \mathcal K(w_1)=w_1,
\]
and
\[
 \mathcal K(w_1,\ldots,w_k)
 =w_k\mathcal K(w_1,\ldots,w_{k-1})
  +\mathcal K(w_1,\ldots,w_{k-2}).
\]
Write \(|\mathbf w|_1=w_1+\cdots+w_k\), with
\(|\varnothing|_1=0\), and let \(\ell(\mathbf w)=k\).  Put
\[
 \begin{aligned}
  \mathcal W_{\rm L}&:=\{\varnothing\}\cup
    \bigcup_{k\ge1}\mathbb N^k,\\
  \mathcal W_{\rm R}&:=\{\varnothing\}\cup
    \{(w_1,\ldots,w_k):k\ge1,\ w_i\ge1,\ w_k\ge2\}.
 \end{aligned}
\]
Thus \(\mathcal W_{\rm R}\) is the set of canonical regular
continued-fraction words, together with the empty word, whereas no terminal
condition is imposed on \(\mathcal W_{\rm L}\).

We next connect \(Z_n(s)\) to the denominator layers used later.  For a
reduced fraction \(p/q\in(0,1)\), write its unique negative continued
fraction as
\[
 \frac pq=
 \cfrac1{e_1-\cfrac1{e_2-\ddots-\cfrac1{e_k}}},
 \qquad e_i\ge2,
\]
and set
\[
 d(p/q):=\sum_{i=1}^k(e_i-1),
 \qquad c(p/q):=2d(p/q)+1.
\]
The negative continuant is defined by
\[
 \begin{gathered}
  \mathcal D(\varnothing)=1,\qquad \mathcal D(e_1)=e_1,\\
  \mathcal D(e_1,\ldots,e_k)
  =e_k\mathcal D(e_1,\ldots,e_{k-1})
   -\mathcal D(e_1,\ldots,e_{k-2}).
 \end{gathered}
\]
Thus \(q=\mathcal D(e_1,\ldots,e_k)\).  For real \(s\), put
\begin{equation}
 b_\ell(s):=
 \sum_{\substack{0<p<q,\ (p,q)=1\\c(p/q)=\ell}}q^{-s}.
 \label{eq:weighted-letter-definition}
\end{equation}
Equivalently,
\begin{equation}
 b_{2d+1}(s)
 =\sum_{k=1}^d
   \sum_{\substack{f_1+\cdots+f_k=d\\f_i\ge1}}
   \mathcal D(f_1+1,\ldots,f_k+1)^{-s},
 \qquad b_{2d}(s)=0.
 \label{eq:weighted-letter-coefficients}
\end{equation}

\begin{proposition}[Exact identification with Dushistova's sum]
\label{prop:dushistova-exact-identification}
For \(n\ge2\) and \(\beta>1\), put
\[
 \mathcal A_n^{\rm D}:=
 \bigcup_{r\ge1}\{(a_1,\ldots,a_r)\in\mathbb N^r:
                         a_r\ge2,\ a_1+\cdots+a_r=n\}
\]
and
\[
 \sigma_\beta(n):=\sum_{\mathbf a\in\mathcal A_n^{\rm D}}
                    \mathcal K(\mathbf a)^{-2\beta}.
\]
Then, term by term,
\begin{equation}
 Z_n(s)=\sigma_{s/2}(n),
 \qquad
 b_{2d+1}(s)=\sigma_{s/2}(d+1)=Z_{d+1}(s)
 \quad(d\ge1,\ s>2).
 \label{eq:dushistova-exact-identification}
\end{equation}
In particular, the field order and renewal indices are related by
\(n=d+1\) and \(c=2d+1=2n-1\).
\end{proposition}

\begin{proof}
The denominator of \([0;a_1,\ldots,a_r]\) is
\(\mathcal K(a_1,\ldots,a_r)\), so the first equality is immediate from
the definitions.  When \(r\) is even, the elementary conversion to the
negative expansion is
\[
 (a_1+1,2^{[a_2-1]},a_3+2,2^{[a_4-1]},\ldots,
   a_{r-1}+2,2^{[a_r-1]}),
\]
where \(2^{[j]}\) denotes a string of \(j\) digits equal to two.  When
\(r>1\) is odd, use the same alternating construction through \(a_{r-1}\),
followed by \(a_r+1\); for \(r=1\), the negative word is \((a_1)\).
Induction in the continued-fraction recursions shows that the paired words
represent the same fraction and have the same denominator.  In all cases,
\[
 d(p/q)=a_1+\cdots+a_r-1.
\]
Hence \(c(p/q)=2d+1\) is equivalent to total regular digit sum \(d+1\),
and the second equality follows term by term.
\end{proof}

\begin{lemma}[Endpoint correction in the left context]
\label{lem:dushistova-endpoint-correction}
For \(s>2\), put
\[
 r_m(s):=\sum_{\substack{\mathbf v\in\mathcal W_{\rm R}\\
                          |\mathbf v|_1=m}}
          \mathcal K(\mathbf v)^{-s},
 \qquad
 \ell_m(s):=\sum_{\substack{\mathbf u\in\mathcal W_{\rm L}\\
                             |\mathbf u|_1=m}}
          \mathcal K(\mathbf u)^{-s}.
\]
Then
\[
 r_0=\ell_0=1,\qquad r_1=0,\qquad \ell_1=1,
 \qquad \ell_m=2r_m\quad(m>1).
\]
Consequently,
\begin{equation}
 \sum_{\mathbf v\in\mathcal W_{\rm R}}\mathcal K(\mathbf v)^{-s}
 =R_s:=\frac{\zeta(s-1)}{\zeta(s)}=\rho_s,
 \qquad
 \sum_{\mathbf u\in\mathcal W_{\rm L}}\mathcal K(\mathbf u)^{-s}
 =2R_s,
 \label{eq:brocot-context-masses}
\end{equation}
and the three left-context cases contribute
\begin{equation}
 \begin{array}{c|c}
  |\mathbf u|_1=0 & R_s\\
  |\mathbf u|_1=1 & R_s\\
  |\mathbf u|_1>1 & 2(R_s-1)R_s
 \end{array}
 \qquad\Longrightarrow\qquad 2R_s^2.
 \label{eq:left-context-endpoint-cases}
\end{equation}
In particular, the conflicting printed and corrected coefficients are
\begin{equation}
 \boxed{\text{Dushistova: }R_s+2R_s^2}
 \qquad\text{and}\qquad
 \boxed{\text{this paper: }2R_s^2}.
 \label{eq:dushistova-conflicting-constants}
\end{equation}
At \(s=\sigma_0\), these are \(10\) and \(8\), respectively.
\end{lemma}

\begin{proof}
Canonical regular words of denominator \(q\) parameterize the
\(\varphi(q)\) reduced fractions in \((0,1)\), where \(\varphi\) is Euler's
totient function.  Including the empty
word and using the totient Dirichlet series gives
\[
 1+\sum_{q\ge2}\frac{\varphi(q)}{q^s}
 =\frac{\zeta(s-1)}{\zeta(s)}=R_s.
\]
For \(m>1\), every canonical word of digit sum \(m\) has exactly two
positive expansions: itself and the word obtained by replacing its final
digit \(a\ge2\) by \(a-1,1\).  This move preserves the digit sum and the
continuant, so \(\ell_m=2r_m\).  At the endpoints, the empty word is the
only word of sum zero, whereas \((1)\) is the only positive word of sum one
and is not canonical.  This proves the identities and
\eqref{eq:left-context-endpoint-cases}.

The precise loss in Dushistova's Lemma~7
\cite[pp.~668--669]{Dushistova2007} is the loss of the restriction \(u>1\).
After separating the \(u=1\) term, the proof replaces the \(u>1\)
left-context sum by twice a canonical convolution denoted \(\Sigma_2\).
Its displayed definition drops \(u>1\), thereby including the empty left
context \(u=0\).  Since \(\Sigma_2\) is doubled, that endpoint is counted
twice rather than once.  With Dushistova's finite-cutoff notation, extending
\(\sigma_\beta(0)=1\) and \(\sigma_\beta(1)=0\), the exact finite identity is
\[
 R_0(w)=\Sigma_1(w)+2\Sigma_2(w)
        -\sum_{v\le w}\sigma_\beta(v).
\]
The subtraction removes exactly one empty-left-context contribution.  Its
limit is \(R_s\), so the printed split has excess exactly \(R_s\).
Retaining the three endpoint classes gives
\(R_s+R_s+2(R_s-1)R_s=2R_s^2\).  This uses the same
empty-continuant, final-digit, and exponent conventions as the published
paper; it is not a normalization change.
\end{proof}

\begin{lemma}[Greedy moderate-block bound]
\label{lem:greedy-moderate-block-bound}
For \(s>2\), define
\[
 \eta_s(h):=
 \sum_{\substack{\mathbf u\in\mathcal W_{\rm L}\\|\mathbf u|_1\ge h}}
 \mathcal K(\mathbf u)^{-s}
\]
and let \(G_s(N,h)\) be the total inverse-continuant weight of the words
\(\mathbf u\in\mathcal W_{\rm L}\) with digit sum at least \(N\) and every
digit smaller than \(h\).  For all sufficiently large \(h\),
\begin{equation}
 G_s(N,h)\le
 \frac{2\rho_s}{1-\eta_s(h)}
 \eta_s(h)^{\lfloor N/(2h)\rfloor-1}.
 \label{eq:moderate-block-bound}
\end{equation}
\end{lemma}

\begin{proof}
Scan a word from the left and close a block at the first digit for which the
block's digit sum reaches \(h\); repeat, leaving the final remainder
unclosed.  This factorization is unique.  Since every digit is smaller than
\(h\), each closed block has digit sum in \([h,2h)\), and a word of digit sum
at least \(N\) has at least \(\lfloor N/(2h)\rfloor-1\) closed blocks.
The concatenation inequality makes the inverse-continuant weight at most the
product of the block weights.  Enlarge each closed-block class to all words
of digit sum at least \(h\), of mass \(\eta_s(h)\), and the remainder class
to all of \(\mathcal W_{\rm L}\), of mass \(2\rho_s\).  Summing the resulting
geometric series proves the bound; \(\eta_s(h)\to0\) by
\eqref{eq:brocot-context-masses}.
\end{proof}

\begin{proposition}[One-large-partial-quotient estimates]
\label{prop:brocot-context-estimates}
Let \(P_n(s)\) be the part of \(Z_n(s)\) contributed by words with a digit
greater than \(n/2\).  For every fixed \(s>2\),
\begin{align}
 n^sP_n(s)&\longrightarrow2\rho_s^2,
 \label{eq:principal-large-digit-limit}\\
 Z_n(s)-P_n(s)&=o(n^{-s}),
 \label{eq:noncondensed-negligible}\\
 Z_n(s)&\sim2\rho_s^2n^{-s}.
 \label{eq:brocot-partition-asymptotic}
\end{align}
More precisely, cutting at the unique digit greater than \(n/2\) is a
bijection with pairs
\((\mathbf u,\mathbf v)\in\mathcal W_{\rm L}\times\mathcal W_{\rm R}\)
such that \(|\mathbf u|_1+|\mathbf v|_1<n/2\), the central digit being
\(a=n-|\mathbf u|_1-|\mathbf v|_1\).  For every fixed pair,
\begin{equation}
 n^s\mathcal K(\mathbf u,a,\mathbf v)^{-s}
 \longrightarrow
 \{\mathcal K(\mathbf u)\mathcal K(\mathbf v)\}^{-s},
 \label{eq:fixed-context-limit}
\end{equation}
and throughout this parameterization,
\begin{equation}
 n^s\mathcal K(\mathbf u,a,\mathbf v)^{-s}
 \le2^s\{\mathcal K(\mathbf u)\mathcal K(\mathbf v)\}^{-s}.
 \label{eq:summable-context-domination}
\end{equation}
\end{proposition}

\begin{proof}
The continuant concatenation formula gives
\begin{equation}
 \mathcal K(\mathbf u\mathbf v)
 \ge\mathcal K(\mathbf u)\mathcal K(\mathbf v),
 \label{eq:continuant-concatenation-lower-bound}
\end{equation}
and multilinearity gives, with the contexts fixed,
\begin{equation}
 \mathcal K(\mathbf u,a,\mathbf v)
 =a\mathcal K(\mathbf u)\mathcal K(\mathbf v)
  +O_{\mathbf u,\mathbf v}(1).
 \label{eq:continuant-large-digit-affine}
\end{equation}
A digit greater than \(n/2\) is unique.  Cutting there gives the asserted
bijection.  If \(\mathbf v\ne\varnothing\), its terminal digit remains at
least two; if \(\mathbf v=\varnothing\), the central digit is greater than
\(n/2\) and is therefore canonical for \(n\ge3\).  Since \(a>n/2\),
\eqref{eq:continuant-concatenation-lower-bound} yields
\eqref{eq:summable-context-domination}, while
\eqref{eq:continuant-large-digit-affine} yields
\eqref{eq:fixed-context-limit}.  The majorant is summable by
\eqref{eq:brocot-context-masses}.  Dominated convergence therefore gives
\eqref{eq:principal-large-digit-limit}.

It remains to bound every word without a digit greater than \(n/2\).  Choose
\[
 \frac{s}{2(s-1)}<\delta<1,
 \qquad h=\lfloor n^\delta\rfloor.
\]
A word having at least two digits not smaller than \(h\) can be cut at its
first two such digits.  Dropping the restrictions on the three intervening
contexts and applying \eqref{eq:continuant-concatenation-lower-bound} bounds
its total weight by
\begin{equation}
 (2\rho_s)^2\rho_s
 \left(\sum_{a\ge h}a^{-s}\right)^2
 =O_s(h^{2-2s})=o(n^{-s}).
 \label{eq:two-large-digits-negligible}
\end{equation}

For words whose digits are all smaller than \(h\),
\Cref{lem:greedy-moderate-block-bound} with \(N=n\) is
superpolynomially small, since \(n/h\to\infty\) and \(\eta_s(h)\to0\).
Every remaining word has exactly one digit \(a\ge h\), with \(a\le n/2\).
Its two surrounding contexts have combined digit sum at least \(n/2\), so
one has digit sum at least \(n/4\).  Cutting at \(a\), bounding the other
context by \(2\rho_s\), and allowing either context to be the long one gives
\[
 4\rho_s\sum_{h\le a\le n/2}a^{-s}G_s(n/4,h).
\]
The first two factors grow at most polynomially, while the last is
superpolynomially small.  Together with
\eqref{eq:two-large-digits-negligible}, this proves
\eqref{eq:noncondensed-negligible}, and hence
\eqref{eq:brocot-partition-asymptotic}.
\end{proof}

\begin{proof}[Proof of \Cref{thm:arithmetic-denominator-layer}]
Combine \Cref{prop:dushistova-exact-identification} with
\eqref{eq:brocot-partition-asymptotic}, and use
\((d+1)^{-s}\sim d^{-s}\).
\end{proof}

\begin{proof}[Proof of \Cref{thm:brocot-condensation}]
For \((\mathbf u,\mathbf v)\in
\mathcal W_{\rm L}\times\mathcal W_{\rm R}\), define the subprobability
mass
\[
 \nu_n(\mathbf u,\mathbf v):=
 \frac{\mathbf1_{\{|\mathbf u|_1+|\mathbf v|_1<n/2\}}}{Z_n(s)}
 \mathcal K(\mathbf u,
 n-|\mathbf u|_1-|\mathbf v|_1,\mathbf v)^{-s}.
\]
Its total mass is \(\mathbb P_{n,s}\{M_n>n/2\}\).  By
\eqref{eq:noncondensed-negligible} and
\eqref{eq:brocot-partition-asymptotic}, this mass tends to one, proving
\eqref{eq:brocot-unique-condensate}.  For each fixed pair,
\eqref{eq:fixed-context-limit} and
\eqref{eq:brocot-partition-asymptotic} give
\[
 \nu_n(\mathbf u,\mathbf v)\longrightarrow
 \nu(\mathbf u,\mathbf v):=
 \frac{\{\mathcal K(\mathbf u)\mathcal K(\mathbf v)\}^{-s}}
      {2\rho_s^2}.
\]
Equation \eqref{eq:brocot-context-masses} shows that \(\nu\) has total mass
one and is the product of the two laws in
\eqref{eq:limiting-context-laws}.

We use the discrete Scheff\'e argument.  Since the masses are nonnegative,
\[
 \sum_{\mathbf u,\mathbf v}|\nu_n-\nu|
 =\sum_{\mathbf u,\mathbf v}\nu_n+1
  -2\sum_{\mathbf u,\mathbf v}\min(\nu_n,\nu).
\]
Fatou's lemma gives
\(\liminf_n\sum\min(\nu_n,\nu)\ge\sum\nu=1\), while these sums are at most
one.  Hence the displayed \(L^1\) distance tends to zero.  The law of
\((U_n,V_n)\) differs from \(\nu_n\) only by placing the complementary
mass at \((\varnothing,\varnothing)\).  That mass tends to zero, proving
\eqref{eq:brocot-context-tv}.

On the event in \eqref{eq:brocot-unique-condensate},
\[
 n-M_n=|U_n|_1+|V_n|_1,
 \qquad
 (J_n-1,r_n-J_n)=(\ell(U_n),\ell(V_n)).
\]
Total variation contracts under a measurable map, and changing a random
variable on an event of probability tending to zero changes its law by at
most that probability.  These observations prove
\eqref{eq:brocot-defect-limit}, \eqref{eq:brocot-defect-law}, and
\eqref{eq:brocot-location-limit}.  Total-variation convergence to the
finite random variable \(D\) makes \(n-M_n\) tight; dividing by \(n\) gives
\(M_n/n\to1\) in probability.

Finally, \eqref{eq:continuant-large-digit-affine} gives, for each fixed
context pair,
\[
 \frac{\mathcal K(\mathbf u,a,\mathbf v)}
      {a\mathcal K(\mathbf u)\mathcal K(\mathbf v)}\longrightarrow1
 \qquad(a\to\infty).
\]
Given \(\varepsilon>0\), total-variation convergence permits a finite set
of context pairs carrying limiting probability at least \(1-\varepsilon\).
On that finite set the last convergence is uniform, and
\(a=n-|\mathbf u|_1-|\mathbf v|_1\to\infty\).  The complement has
limsup probability at most \(\varepsilon\), apart from the vanishing
noncondensed event.  Letting \(\varepsilon\downarrow0\) proves
\eqref{eq:brocot-denominator-ratio}.
\end{proof}

\section{The Fibonacci generator renewal}\label{sec:fibonacci-renewal}

Let \(F_0=0,F_1=1\), and \(F_{j+2}=F_{j+1}+F_j\).  The classical Fibonacci
partition function is
\begin{equation}
 R(N):=\card{\left\{(\varepsilon_j)_{j\ge0}\in\{0,1\}^{(\mathbb N)}:
 N=\sum_{j\ge0}\varepsilon_jF_{j+2}\right\}},
 \qquad N\ge0.
 \label{eq:fibonacci-partition-function}
\end{equation}
Thus the available parts are \(1,2,3,5,\ldots\), each used at most once.
For \(m\ge1\), put
\begin{equation}
 \mathcal I_m=[F_{m+1}-1,F_{m+2}-1)\cap\ZZ,
 \qquad
 Z_m^R(t):=\sum_{N\in\mathcal I_m}R(N)^t.
 \label{eq:standard-fibonacci-layer}
\end{equation}

We use the generator coordinates of Weinstein
\cite[Secs.~3--5]{Weinstein2016}.  His Fibonacci sequence
\(f_0=f_1=1\) satisfies \(f_j=F_{j+1}\), so his half-open layer
\([f_m-1,f_{m+1}-1)\) is exactly \(\mathcal I_m\), and his partition
function is \(R\) in \eqref{eq:fibonacci-partition-function}.  We record the
part of the parameterization needed for weighted sums.

Let
\[
 \mathcal A=\{p/q:0<p<q,\ (p,q)=1\}.
\]
For \(p/q\in\mathcal A\), write its unique negative continued fraction as
\[
 \frac pq=
 \cfrac1{e_1-\cfrac1{e_2-\ddots-\cfrac1{e_k}}},
 \qquad e_i\ge2,
\]
and recall from \Cref{sec:brocot-condensation} that
\[
 d(p/q)=\sum_{i=1}^k(e_i-1),
 \qquad c(p/q)=2d(p/q)+1.
\]
For real \(s\), define
\begin{equation}
 b_\ell(s)=\sum_{\substack{p/q\in\mathcal A\\c(p/q)=\ell}}q^{-s},
 \qquad
 B_s(z)=\sum_{\ell\ge1}b_\ell(s)z^\ell.
 \label{eq:renewal-letter-series}
\end{equation}
Only odd costs occur.  The smallest costs are \(3\) and \(5\), represented
by \(1/2\) and \(1/3\).

Put
\begin{equation}
 u_0(s)=1,
 \qquad
 u_j(s)=\sum_{\ell=1}^jb_\ell(s)u_{j-\ell}(s)\quad(j\ge1).
 \label{eq:weighted-renewal-coefficients}
\end{equation}
Equivalently, in formal power series,
\begin{equation}
 U_s(z):=\sum_{j\ge0}u_j(s)z^j=\frac1{1-B_s(z)}.
 \label{eq:weighted-renewal-generating-function}
\end{equation}

\begin{proposition}[Weighted generator renewal]
\label{prop:weighted-generator-renewal}
For every real \(s\) and \(m\ge1\),
\begin{equation}
 Z_m^R(-s)
 =1+u_m(s)+2\sum_{j=1}^{m-1}u_j(s)
 =2\sum_{j=0}^{m-1}u_j(s)+u_m(s)-1.
 \label{eq:standard-layer-renewal}
\end{equation}
Moreover, \(u_j(s)\) is exactly the total inverse-\(s\) weight of the
Weinstein generators of cost \(j\).
\end{proposition}

\begin{proof}
Under Weinstein's free-monoid bijection, a generating number \(g\)
corresponds uniquely to a nonempty word
\[
 w=\frac{p_1}{q_1}\times\cdots\times\frac{p_r}{q_r}
 \qquad(p_i/q_i\in\mathcal A).
\]
Proposition~3.3 and formula~(16) of
\cite{Weinstein2016}, translated through the preceding Fibonacci convention,
give
\begin{equation}
 R(g)=\prod_{i=1}^rq_i,
 \qquad
 L(g)=\sum_{i=1}^rc(p_i/q_i),
 \label{eq:generator-weight-cost}
\end{equation}
where \(L(g)\) is the largest Zeckendorf index of \(g\).  Unique
factorization in the free monoid and \eqref{eq:renewal-letter-series} show
that the total weight \(R(g)^{-s}\) of generators with \(L(g)=j\) is the
coefficient \(u_j(s)\).  They also prove
\eqref{eq:weighted-renewal-generating-function}.

It remains to count one generating orbit in a fixed half-open Fibonacci
layer.  Weinstein's normal form consists of the four branches
\([a]g,[a]\tau g,\sigma([a]g),\sigma([a]\tau g)\), with \(a\ge0\).  His
formulas~(17)--(19) give the corresponding largest indices
\[
 \begin{array}{c|c}
 \text{branch}&\text{largest index}\ \\ \hline
 {}[a]g&L(g)+2a\\
 {}[a]\tau g&L(g)+1+2a\\
 \sigma({}[a]g)&L(g)+1+2a\\
 \sigma({}[a]\tau g)&L(g)+2+2a.
 \end{array}
\]
The two middle entries are distinct because the action is free.  Hence an
orbit of generator cost \(j=L(g)\) contributes one point to
\(\mathcal I_j\), two points to every \(\mathcal I_m\) with \(m>j\), and
no point to earlier layers.  Distinct generators have disjoint orbits.
There is in addition exactly one partition of multiplicity \(1\) in each
\(\mathcal I_m\), namely its left endpoint.  Summing the orbit weights
therefore gives the first expression in \eqref{eq:standard-layer-renewal};
the second is algebraic rearrangement.
\end{proof}

The critical normalization follows directly from the totient Dirichlet
series.  For \(s>2\), every denominator \(q\ge2\) occurs with its
\(\varphi(q)\) reduced numerators, so Tonelli's theorem gives
\begin{equation}
 B_s(1)=\sum_{q\ge2}\frac{\varphi(q)}{q^s}
 =\frac{\zeta(s-1)}{\zeta(s)}-1.
 \label{eq:weighted-renewal-at-one}
\end{equation}
The right side decreases continuously from infinity to zero as \(s\) runs
from \(2\) to infinity.  Thus there is a unique \(\sigma_0>2\) with
\begin{equation}
 \frac{\zeta(\sigma_0-1)}{\zeta(\sigma_0)}=2,
 \qquad B_{\sigma_0}(1)=1.
 \label{eq:critical-renewal-normalization}
\end{equation}

At this exponent, \(b_j(\sigma_0)\) is an aperiodic probability law: its
support contains \(3\) and \(5\).  Its mean
\begin{equation}
 \mu_C:=\sum_{q\ge2}\sum_{\substack{1\le p<q\\(p,q)=1}}
 c(p/q)q^{-\sigma_0}
 \label{eq:critical-cost-mean}
\end{equation}
is finite.  Indeed, conversion to the canonical regular continued fraction
\([0;a_1,\ldots,a_r]\) gives
\(c(p/q)=2(a_1+\cdots+a_r)-1\).  The fixed-denominator mean theorem of
Panov and Liehl \cite{Panov1982,Liehl1983} bounds the sum of
\(a_1+\cdots+a_r\) over reduced numerators by
\(O(\varphi(q)(\log q)^2)\).  Therefore
\[
 \mu_C\ll\sum_{q\ge2}q^{1-\sigma_0}(\log q)^2<\infty.
\]
This verifies the finite-mean hypothesis used in the critical renewal
argument below.

\section{The arithmetic critical law}\label{sec:arithmetic-critical}

\begin{theorem}[Arithmetic critical stable renewal]
\label{thm:arithmetic-critical-stable-renewal}
Put \(\alpha=\sigma_0-1\in(1,2)\), and under the critical letter law put
\[
 \Pr\{(C,H)=(c(p/q),\log q)\}=q^{-\sigma_0}
 \qquad(p/q\in\mathcal A).
\]
Then \(\Pr\{C=2d+1\}=b_{2d+1}(\sigma_0)\), and the constants
\begin{equation}
 \begin{aligned}
 b_C&:=
 \sum_{\mathbf u\in\mathcal W_{\rm L}}
 \sum_{\mathbf v\in\mathcal W_{\rm R}}
 \{\mathcal K(\mathbf u)\mathcal K(\mathbf v)\}^{-\sigma_0}
 =2\left(\frac{\zeta(\sigma_0-1)}{\zeta(\sigma_0)}\right)^2=8,\\
 K_C&:=\frac{2^{\sigma_0-1}}{\sigma_0-1}b_C
 =\frac{2^{\sigma_0+2}}{\sigma_0-1}
 \end{aligned}
 \label{eq:critical-tail-constants}
\end{equation}
are positive.  The double series defining \(b_C\) is absolutely convergent,
and
\begin{align}
 b_{2d+1}(\sigma_0)&\sim b_Cd^{-\sigma_0},
 &&d\to\infty,
 \label{eq:critical-letter-local-tail}\\
 \Pr\{C>x\}
 &=\sum_{q\ge2}q^{-\sigma_0}
   \card{\{1\le p<q:(p,q)=1,\ c(p/q)>x\}}\notag\\
 &\sim K_Cx^{1-\sigma_0},
 &&x\to\infty.
 \label{eq:critical-letter-tail}
\end{align}
If \(u_j=u_j(\sigma_0)\) is the renewal sequence in
\eqref{eq:weighted-renewal-coefficients}, then
\begin{equation}
 u_j=\frac1{\mu_C}
 +\frac{K_C}{\mu_C^2(\sigma_0-2)}j^{2-\sigma_0}
 +o(j^{2-\sigma_0}).
 \label{eq:critical-renewal-second-order}
\end{equation}
Consequently,
\begin{equation}
 Z_m^R(-\sigma_0)
 =\frac{2m}{\mu_C}
 +\frac{2K_C}{\mu_C^2(\sigma_0-2)(3-\sigma_0)}
   m^{3-\sigma_0}+o(m^{3-\sigma_0}).
 \label{eq:standard-layer-critical-second-order}
\end{equation}
Finally, \(C\) belongs to the domain of attraction of a spectrally positive
\(\alpha\)-stable law.  More precisely, if \(a_n>0\) is chosen so that
\(n\Pr\{C>a_n\}\to1\), then for independent copies \(C_1,C_2,\ldots\),
\begin{equation}
 \frac{C_1+\cdots+C_n-n\mu_C}{a_n}
 \xrightarrow{\mathrm d}\mathcal S_\alpha,
 \label{eq:critical-cost-stable-domain}
\end{equation}
where the centered spectrally positive law \(\mathcal S_\alpha\) is fixed by
\begin{equation}
 \log\EE e^{it\mathcal S_\alpha}
 =\int_0^\infty
   (e^{itx}-1-itx)\alpha x^{-\alpha-1}\,dx.
 \label{eq:centered-stable-characteristic-exponent}
\end{equation}
\end{theorem}

\begin{proof}
Equation \eqref{eq:critical-renewal-normalization} makes the displayed
letter law a probability measure.  The local estimate
\eqref{eq:critical-letter-local-tail} is
\Cref{thm:arithmetic-denominator-layer} at \(s=\sigma_0\), and
\eqref{eq:brocot-context-masses} evaluates the absolutely convergent context
sum in \eqref{eq:critical-tail-constants}.  Direct summation gives
\[
 \Pr\{C>x\}
 =\sum_{d>(x-1)/2}b_{2d+1}(\sigma_0)
 \sim\frac{b_C}{\sigma_0-1}
      \left(\frac{x}{2}\right)^{1-\sigma_0},
\]
which is \eqref{eq:critical-letter-tail}.  Thus regular variation is a
consequence of the arithmetic continuant estimate.

Write \(\overline F(j)=\Pr\{C>j\}\).  The increments are nonnegative
integers, their mean \(\mu_C\) is finite by
\eqref{eq:critical-cost-mean}, and their span is one because the fractions
\(1/2\) and \(1/3\) have costs \(3\) and \(5\).  The hypotheses of the
second-order lattice renewal theorem of Omey and Van Gulck are therefore
satisfied.  Their asymptotic
\cite[Sec.~3.2]{OmeyVanGulck2015} is
\begin{equation}
 u_j-\frac1\mu
 \sim\frac{j\overline F(j)}{\mu^2(\alpha-1)}
 \label{eq:second-order-renewal-input}
\end{equation}
for an aperiodic nonnegative integer law of finite mean and tail regularly
varying with index \(-\alpha\), \(\alpha>1\).  Taking
\(\alpha=\sigma_0-1\) proves
\eqref{eq:critical-renewal-second-order}.

Summing this expansion and using
\[
 \sum_{j=1}^{m-1}j^{2-\sigma_0}
 \sim\frac{m^{3-\sigma_0}}{3-\sigma_0}
\]
in the exact identity \eqref{eq:standard-layer-renewal} proves
\eqref{eq:standard-layer-critical-second-order}; the endpoint terms are
\(O(1)=o(m^{3-\sigma_0})\).

Finally, \eqref{eq:critical-letter-tail}, finite mean, and the absence of a
negative tail are the one-sided domain-of-attraction conditions for an
\(\alpha\)-stable law with the normalization and centering in
\eqref{eq:critical-cost-stable-domain} and
\eqref{eq:centered-stable-characteristic-exponent}; see
\cite[Ch.~XVII]{Feller1971}.
\end{proof}

\section{Conclusion}\label{sec:conclusion}

The total-variation theorem describes the sampled continued fraction beyond
its scalar normalization.  It gives simultaneous limits for the two
contexts, the condensate defect and location, and the denominator
factorization.  A balanced cut and the exact boundary terms in continuant
concatenation give the sharp \(n^{-1}\) total-variation rate and its limiting
constant.  The summable two-context estimate is the arithmetic step that
makes these joint conclusions possible for the nonmultiplicative continuant
weight.

Its scalar denominator-layer asymptotic supplies the critical cost tail in
Weinstein's generator coordinates.  The exact orbit count turns the
classical Fibonacci layer into a renewal sum, and the tail then produces the
unbounded sublinear term of order \(m^{3-\sigma_0}\).  The same tail yields
the stable domain of attraction for generator costs.  Combined with the
orbit count, it also gives the joint critical Gibbs law: generator cost is
asymptotically uniform in the layer, with a stable fluctuation in the number
of generator letters.

The remaining question is whether analogous critical renewal corrections
hold for other nonmultiplicative continued-fraction families.  Such an
extension would require an arithmetic generator identity and a local tail
estimate adapted to the family.

\begin{thebibliography}{99}

\bibitem{ArmendarizLoulakis2011}
I.~Armend\'ariz and M.~Loulakis,
Conditional distribution of heavy tailed random variables on large
deviations of their sum,
\textit{Stochastic Process. Appl.} \textbf{121} (2011), no.~5, 1138--1147,
doi:10.1016/j.spa.2011.01.011.

\bibitem{Berstel2001}
J.~Berstel,
An exercise on Fibonacci representations,
\textit{RAIRO Theor. Inform. Appl.} \textbf{35} (2001), no.~6, 491--498,
doi:10.1051/ita:2001127.

\bibitem{Carlitz1968}
L.~Carlitz,
Fibonacci representations,
\textit{Fibonacci Quart.} \textbf{6} (1968), no.~4, 193--220.

\bibitem{Dushistova2007}
A.~A.~Dushistova,
Partitioning of the interval \([0,1]\) induced by the Brocot sequences,
\textit{Sb. Math.} \textbf{198} (2007), no.~5, 661--690,
doi:10.1070/SM2007v198n05ABEH003854.

\bibitem{Feller1971}
W.~Feller,
\textit{An Introduction to Probability Theory and Its Applications}, vol.~II,
2nd ed., Wiley, New York, 1971.

\bibitem{Liehl1983}
B.~Liehl,
\"Uber die Teilnenner endlicher Kettenbr\"uche,
\textit{Arch. Math. (Basel)} \textbf{40} (1983), no.~2, 139--147,
doi:10.1007/BF01192762.

\bibitem{MoshchevitinZhigljavsky2004}
N.~Moshchevitin and A.~Zhigljavsky,
Entropies of the partitions of the unit interval generated by the Farey tree,
\textit{Acta Arith.} \textbf{115} (2004), no.~1, 47--58,
doi:10.4064/aa115-1-4.

\bibitem{OmeyVanGulck2015}
E.~Omey and S.~Van Gulck,
Intuitive approximations in discrete renewal theory, Part 1: Regularly
varying case,
\textit{Statist. Probab. Lett.} \textbf{104} (2015), 68--74,
doi:10.1016/j.spl.2015.05.002.

\bibitem{Panov1982}
A.~A.~Panov,
The mean for a sum of elements in a class of finite continued fractions,
\textit{Math. Notes} \textbf{32} (1982), no.~5, 781--785,
doi:10.1007/BF01358470.

\bibitem{Stufler2020}
B.~Stufler,
Unlabelled Gibbs partitions,
\textit{Combin. Probab. Comput.} \textbf{29} (2020), no.~2, 293--309,
doi:10.1017/S0963548319000336.

\bibitem{Weinstein2016}
F.~V.~Weinstein,
Notes on Fibonacci partitions,
\textit{Exp. Math.} \textbf{25} (2016), no.~4, 482--499,
doi:10.1080/10586458.2015.1118416, arXiv:math/0307150.

\bibitem{Zeckendorf1972}
E.~Zeckendorf,
Repr\'esentation des nombres naturels par une somme de nombres de Fibonacci ou
de nombres de Lucas,
\textit{Bull. Soc. Roy. Sci. Li\`ege} \textbf{41} (1972), 179--182.

\end{thebibliography}


\end{document}
